\section{Temporal results}
\label{sec:temporal}

The original is cross-sectional by construction, and its conclusion names the extension it
does not attempt: ``The cross-sectional nature of our data limits causal inferences,
highlighting the need for longitudinal studies to examine how changes in corporate
characteristics and regulatory environments affect ESGSI over time'' \citep[p.~10]{lagasio2024}.
This section supplies that series on our panel.

It is reported last among the empirical sections, and it is subordinate to what precedes it.
\S\ref{sec:measurement} establishes that a TF-IDF substance score measures the \emph{mix} of
ESG terms rather than their amount, and \S\ref{sec:artefact} that the original's headline
conclusion follows from an undisclosed score-normalisation choice. Having argued that the
instrument is unsound under one admissible specification, we cannot then present a time series
of it as the paper's contribution. What follows is the payload the corrected instrument
delivers, conditional on the choices justified in \S\ref{sec:validity} and bounded by the
sensitivity documented in \S\ref{sec:robustness}. The central claim remains
\S\ref{sec:artefact}'s. By the end of the section the reason for the subordination will be
sharper than a matter of ordering: \S\ref{sec:temporal-mandate} identifies a mechanism under
which the series would move exactly as it does with no change in firm conduct whatever, and
that mechanism is a property of the index architecture rather than of our panel.

The series is computed under $\ESGSI(\susdens)$, the specification \S\ref{sec:validity} selects
on convergent validity against $\QUANT$. We report the base index only: the extended index
carries two arbitrary weights, and \S\ref{sec:limitations} sets out why its temporal slope is
contaminated rather than cleaner.

\subsection{The series}

Table~\ref{tab:temporal-series} gives the yearly mean. The panel is complete: 49 firms, seven
years, 343 documents, exactly 49 in every year, with no entry, exit or backfill. The mean falls
from $0.563$ in 2018 to $-0.674$ in 2024 and is strictly decreasing in every year of the
window. An OLS regression of the document-level index on calendar year ($n = 343$) gives a
slope of $-0.205$ per year with $p = 1.7\times10^{-8}$.

\begin{table}[htbp]
\centering
\caption{Mean $\ESGSI$ by year, under the adopted substance specification and under the
$L_2$-normalised alternative. Balanced panel, 49 documents per year. Values are read from the
pipeline output, which standardises with the population standard deviation.}
\label{tab:temporal-series}
\small
\begin{tabular}{@{}lrrrrrrr@{}}
\toprule
 & 2018 & 2019 & 2020 & 2021 & 2022 & 2023 & 2024 \\
\midrule
$\ESGSI(\susdens)$ (adopted) & $0.563$ & $0.471$ & $0.161$ & $0.007$ & $-0.121$ & $-0.407$ & $-0.674$ \\
$\ESGSI(\suslag)$            & $0.561$ & $0.453$ & $0.106$ & $-0.168$ & $-0.274$ & $-0.371$ & $-0.307$ \\
Documents & 49 & 49 & 49 & 49 & 49 & 49 & 49 \\
\bottomrule
\end{tabular}
\end{table}

Three qualifications attach to that regression before anything is read off it.

The slope is a summary, not a model of the path: the year-on-year decrements are uneven, and
nothing here identifies a constant annual rate. The crossing from positive to negative between
2021 and 2022 is likewise not a calibrated event. Both components are score-normalised across
the pooled sample, so the index has mean zero over the seven years by construction; the
location of the zero crossing is a property of the estimation window, not of the disclosure.
The slope is invariant to that centring; the levels are not.

The $p$-values reported throughout this section, here and in
Table~\ref{tab:temporal-subsample}, come from OLS on document-level observations and treat the
343 documents as independent. They are not. The panel is 49 firms observed seven times each,
and both firm-level heterogeneity and within-firm serial correlation are present by
construction, so the effective number of independent units is nearer 49 than 343 and the
conventional standard errors are anti-conservative. We do not report clustered or
panel-corrected inference. The magnitudes that survive that objection are not the $p$-values
but the facts they summarise: seven consecutive yearly means falling in sequence, a decline
of comparable sign and larger magnitude reappearing in a disjoint single-country subsample,
and driver quantities that are directly counted rather than estimated
(Table~\ref{tab:temporal-drivers}). Readers who discount the exponents should discount them;
the section's argument does not rest on any of them being exact.

\subsection{What moves, and what does not}

$\ESGSI = \Z(\SEN) - \Z(\SUS)$, so a decline is either tone falling or substance rising. The
displacement sits on the substance side. Table~\ref{tab:temporal-drivers} decomposes the
endpoints.

\begin{table}[htbp]
\centering
\caption{Driver decomposition, 2018 against 2024. Per-report means.}
\label{tab:temporal-drivers}
\small
\begin{tabular}{@{}lrrr@{}}
\toprule
Quantity & 2018 & 2024 & Change \\
\midrule
Extracted ESG text per report (lemmatised tokens) & 27{,}906 & 63{,}654 & $+128.1\%$ \\
ESG mentions per report & 1{,}382.7 & 4{,}832.0 & $+249.5\%$ \\
$\susdens$ (mentions per 100 words) & 4.87 & 7.45 & $+52.8\%$ \\
$\BREADTH$ (effective distinct terms) & 28.99 & 41.73 & $+43.9\%$ \\
Extracted zones per report & 178.5 & 275.2 & $+54.2\%$ \\
\bottomrule
\end{tabular}
\end{table}

Three quantities grow together. The volume of text the extraction stage identifies as
ESG-relevant more than doubles; within that text, ESG mentions grow faster still, so density
rises by half again against the larger denominator; and the mentions spread over more of the
vocabulary, raising the effective number of distinct terms by 44\%. Reports are not merely
longer. They are lexically denser in ESG terms and broader in coverage.

$\SEN$ shows no detectable trend. Regressed on calendar year over the same panel, its slope is
$+0.0016$ per year with $p = 0.666$: a point estimate two orders of magnitude below the index
slope, and of the wrong sign to contribute to the decline, since rising tone raises $\ESGSI$.
Two things follow and a third does not. It follows that the movement in the index is located on
the substance side, and that the decline is not produced by firms writing in a less positive
register. It does not follow that tone is uninformative, or that $\SEN$ measures nothing: a
failure to reject at $p = 0.666$ is an absence of detectable movement, not a demonstrated null,
and $\SEN$ retains cross-sectional variance that enters every index value in the panel. We
claim the location of the movement, not the irrelevance of the component.

The null does dispose of one specific objection, and only that one. A dictionary-based index
applied longitudinally invites the charge that a fixed word list scores 2024 language
differently from 2018 language, so that an apparent trend is dictionary drift rather than
disclosure change. For the \emph{sentiment} dictionary that charge would bind only if the
movement were in $\SEN$, and it is not. It does not extend to the ESG vocabulary, which is
where the load is carried and where the equivalent objection is live rather than closed:
\S\ref{sec:temporal-mandate} takes it up, and does not resolve it.

What the series does not establish is that the added substance is real activity. The index
registers that firms wrote more, and more diversely, about ESG topics per unit of text. Whether
that text corresponds to conduct is outside what any text-derived index can determine.

\subsection{Alternative explanations}

Table~\ref{tab:temporal-threats} lists the explanations a downward trend in this index admits
besides a change in disclosure, the check available against each, and what that check settles.
Two rows are not settled: the first of the two is bounded but not removed, the second is
neither, and both are the subject of \S\ref{sec:temporal-mandate}.

\begin{table}[htbp]
\centering
\caption{Alternative explanations for the decline, and the status of each after the checks
this design supports.}
\label{tab:temporal-threats}
\small
\begin{tabular}{@{}p{4.5cm}p{5.6cm}p{3.4cm}@{}}
\toprule
Explanation & Check & Status \\
\midrule
Country and document-type mix drift across years
& Balanced panel: 49 documents and a fixed per-country count every year; 33--36 annual
reports, 12--15 URDs, 1 sustainability report per year (\S\ref{sec:data})
& Excluded by construction \\
The country/type confound generates the slope
& Germany: 98 documents, one country, one document type, $-0.301$/yr
& Excluded within that subsample \\
Reports simply got longer
& $\susdens$ is a ratio; dilution lowers it (\S\ref{sec:measurement})
& Excluded; biases against the finding \\
Extraction admits more text over time
& Marginal zones enter at the density threshold, below the mean; extraction-stage density
rises across the window
& Excluded on the extensive margin \\
ESG language concentrated rather than increased
& None: no whole-document density is computed by year
& Not testable in this design \\
Sentiment dictionary scores 2024 language differently
& $\SEN$ slope $+0.0016$/yr, $p = 0.666$
& No movement to explain \\
ESG vocabulary contains entries whose referents postdate 2018
& Recomputed with the 34-term reporting-frameworks block deleted: slope $-0.155$/yr,
$p = 2.5\times10^{-5}$ (\S\ref{sec:temporal-mandate})
& Bounded, not excluded \\
Mandated disclosure raises term counts without a change in conduct
& None: exposure is common to all firms and collinear with calendar time
& Not excluded (\S\ref{sec:temporal-mandate}) \\
\bottomrule
\end{tabular}
\end{table}

\paragraph{Composition.} The panel is exactly balanced. The same 49 firms appear in all seven
years, every year contributes exactly 49 documents, and every country's own count is identical
in every year (\S\ref{sec:data}). The country mix cannot drift because there is nothing for it
to drift into. Document-type composition is likewise stable: every year contains between 33 and
36 annual reports, between 12 and 15 URDs and exactly one standalone sustainability report, a
year-to-year variation of at most three documents in 49 and not directional.

Balance is bought at a price that belongs here rather than in a later section. Requiring one
report per firm in each of seven consecutive years selects firms that reported continuously
across the window and survived it. The trend is therefore a within-firm trend among continuous
reporters, and the sample it generalises to is that population, not European issuers at large.
This restricts external validity; it does not generate an internal trend, since the selection
is applied identically to every year.

\paragraph{Length.} A report that grows without adding ESG content does not raise $\susdens$; it
lowers it. \S\ref{sec:measurement} demonstrates the point directly: appending 56,000 non-ESG
filler words to a 22,576-word document, with ESG mentions held at 695, drops density from
$3.0785$ to $0.8845$ per 100 words, while the $L_2$-normalised score does not move at all. The
denominator is doing exactly what a density denominator should. Length growth on its own is
therefore not a rival explanation for the decline in $\ESGSI$ but a force working against it,
and the observed rise in density is net of whatever dilution the growth in report length
produced.

\paragraph{Extraction yield.} $\susdens$ counts vocabulary hits over the lemmatised tokens of
the \emph{extracted} zones, not of the whole document, and the extraction stage admits more
text in later years --- zones per report rise from 178.5 to 275.2. If the measure were
mechanically tracking that yield, the growth would be an artefact of the pipeline. It is not,
and the direction of the selection is why. Zones are built around paragraphs that clear a
keyword-density threshold of 1.0 per 100 words, with one paragraph of context admitted on
either side (\S\ref{sec:extraction}). Text admitted at the extensive margin therefore enters at
or near that threshold, well below the corpus mean, and the context windows import adjacent
text with no density requirement at all: adding zones dilutes. The extraction-stage keyword
density --- the \S\ref{sec:extraction} quantity, computed with the full regex over raw text,
and not to be read as $\susdens$ --- nevertheless rises monotonically across the window
rather than falling, and it does so under both readings of that quantity, the pooled ratio
and the mean of per-report values. The intensive margin dominates the extensive one, which is
the opposite of what a yield artefact would produce.

One residual is not disposed of. Extraction and measurement share a vocabulary
(\S\ref{sec:extraction}), so density is measured on text selected for being ESG-dense. A firm
that held its total ESG content fixed but consolidated it into denser, more clearly delimited
passages --- a dedicated non-financial statement in place of remarks scattered through the
report --- would raise the density of the extracted zones without raising the density of the
document. No whole-document density is computed by year, so we cannot separate that from an
increase in content, and the yield of the extraction stage relative to the full PDF was
measured on a 20-document sample rather than year by year. We flag this rather than dismiss it.

\paragraph{The single-country subsample.} The strongest of the composition checks is that the
trend holds inside one country with one document type. Germany contributes 98 documents ---
14 firms over seven years --- and all 98 are annual reports. Within that subsample the slope is
$-0.301$ per year with $p = 1.2\times10^{-6}$, steeper than the pooled estimate rather than
attenuated (Table~\ref{tab:temporal-subsample}).

\begin{table}[htbp]
\centering
\caption{Trend in $\ESGSI(\susdens)$, pooled panel and single-country subsample. OLS of the
document-level index on calendar year. Both $p$-values assume independent observations, which
49 firms observed seven times each are not.}
\label{tab:temporal-subsample}
\small
\begin{tabular}{@{}lrrrl@{}}
\toprule
Sample & $n$ & Slope/yr & $p$ & Composition \\
\midrule
Full panel & 343 & $-0.205$ & $1.7\times10^{-8}$ & 7 countries, 3 document types \\
Germany only & 98 & $-0.301$ & $1.2\times10^{-6}$ & 1 country, 1 document type \\
\bottomrule
\end{tabular}
\end{table}

What this check removes, and what it does not, must be stated exactly.
\S\ref{sec:data} identifies a near-total confound between country and document type and treats
any cross-sectional comparison along either dimension as uninterpretable. Inside the German
subsample neither factor varies, so neither can be generating the slope \emph{there}. That is
the whole of the inference. It does not make the German subsample representative, and it does
not transfer to the pooled estimate by extension: it shows that a trend of this sign and
magnitude exists in a slice where the confound is absent, not that the confound is absent from
the pooled slope. The German slope is in fact steeper than the pooled one, which is itself
evidence that the rate is heterogeneous across the panel, so $-0.301$ is not to be read as a
European rate. The subsample is 98 documents from 14 firms, so the dependence between
observations noted for the pooled regression binds harder here, not less. And it is silent by
construction about
anything acting on all seven countries at once --- which is the next subsection's problem, and
the one this check is least able to address.

\subsection{The objection this design cannot answer}
\label{sec:temporal-mandate}

The window 2018--2024 is the interval over which European non-financial reporting moved from
largely voluntary to substantially prescribed. A firm that adds a taxonomy-alignment table, a
framework index and a standards-referenced non-financial statement adds ESG vocabulary to its
report in quantity, without necessarily changing anything it does. Our numerator counts words.
It cannot distinguish the two, and the objection is the strongest available against reading
this series as evidence about firms.

The mechanism is not hypothetical in our instrument; it is visible in the vocabulary. The
154-term list contains entries --- \emph{csrd}, \emph{esrs}, \emph{issb}, \emph{sfdr},
\emph{tcfd}, and four members of the \emph{taxonomy} family --- that name regimes, standards
and disclosure categories which either did not exist or were not in force at the start of the
window. Their near-absence from 2018 reports is a matter of chronology, not of disclosure
behaviour. A fixed word list applied across a period in which the referents of some of its
entries came into being will register an increase that is to that extent definitional. This is
the ESG-vocabulary counterpart of the dictionary-drift objection that the $\SEN$ null disposes
of on the sentiment side. Here it is not disposed of. The definitional part of it can be
bounded, and below we bound it; the objection as a whole survives that bound.

\S\ref{sec:validity} supplies one measurement of the drift, and it points the wrong way for
us. Fourteen vocabulary entries overlap the \QUANT{} regex --- eleven in the
\texttt{frameworks} group, three in \texttt{units} --- and carry 7.50\% of all $\SUS$
occurrences pooled. That share is not stationary: the \texttt{frameworks} group rises from
10.15\% of all \QUANT{} hits in 2018 to 21.87\% in 2024. Framework and standard references
roughly double as a share of quantified disclosure across exactly the window in which the
series declines. The drift is real, it is measured, and it is in the direction the objection
predicts.

Two counterweights are available, and neither disposes of the confound; the first bounds the
part of it that is purely definitional, the second is an argument only.

The first is to delete the offending vocabulary and recompute. The \emph{reporting frameworks
and ratings} block of
\S\ref{sec:lexicon} is the block in which regime names live, and it is where every entry
above sits. Dropping it in its entirety --- all 34 of its terms, taking the vocabulary from
154 to 120 --- removes far more than the chronologically suspect entries, since the block also
contains framework vocabulary that predates the window. It is deliberately over-inclusive: we
want an upper bound on what the definitional component can be worth, so we discard the whole
block rather than adjudicate entry by entry. Table~\ref{tab:temporal-frameworks} reports the
result. The rise in density falls from $+52.8\%$ to roughly $+43\%$, and the index slope from
$-0.205$ to $-0.155$ per year --- about three quarters of its magnitude --- remaining
significant at any conventional level.

\begin{table}[htbp]
\centering
\caption{Temporal trend with the \emph{reporting frameworks and ratings} block removed from
the ESG vocabulary (154 $\to$ 120 terms). Density is the mean of per-report $\susdens$ values.
$p$-values carry the same independence caveat as everywhere else in this section.}
\label{tab:temporal-frameworks}
\small
\begin{tabular}{@{}lrr@{}}
\toprule
 & Full vocabulary (154) & Frameworks excluded (120) \\
\midrule
Mean density, 2018 $\to$ 2024 & $4.87 \to 7.45$ & $4.43 \to 6.35$ \\
$\ESGSI(\susdens)$ slope & $-0.205$/yr & $-0.155$/yr \\
$p$ & $1.7\times10^{-8}$ & $2.5\times10^{-5}$ \\
\bottomrule
\end{tabular}
\end{table}

What that check bounds and what it leaves untouched must be stated exactly, because it is
easy to overread. It shows that the trend is not an artefact of counting regime names: strip
every one of them, and three quarters of the slope is still there. It does \emph{not} dispose
of the mandate confound, and we do not offer it as doing so. A reporting regime elicits
substantive disclosure as well as its own acronyms --- emissions-scope terms, net-zero
language and the greater part of the environmental vocabulary are prescribed by the same
regimes without naming them --- and that portion of the effect is not separable from conduct
by any vocabulary edit, because the words involved are exactly the words a firm that had
genuinely changed its conduct would also use. The 34-term deletion bounds the purely
definitional component of the confound. The rest of the confound is untouched, and it is the
larger part.

The second counterweight is breadth. A compliance template inserts a bounded set of named items, which
predicts more mentions concentrated on few terms. What we observe is broadening: the distinct
vocabulary entries present per report rise from 68.65 to 105.45 of 154, and $\BREADTH$ from
28.99 to 41.73. That is more than a handful of new acronyms can produce. It is not decisive
either, because the mandates in question prescribe reporting across environmental, social and
governance topics simultaneously, so a mandate effect can broaden coverage as well as deepen
it. The counterweight restates the confound at a higher level rather than resolving it.

No design in this paper can do better than that bound, and it is worth being precise about
why the residue is not reachable. A reporting mandate is a common time shock. It applies to all 49 firms in the same years, so its exposure
is collinear with calendar year, which is the regressor. The balanced panel, the document-type
check and the German subsample all remove \emph{cross-sectional} composition; none of them
removes anything that moves with time. The German result is the weakest defence available here
rather than the strongest, since German issuers face the same European regimes as the rest of
the panel: holding country and document type fixed does nothing to hold the regulatory calendar
fixed. Identification would require variation in exposure --- firms inside and outside scope,
or staggered application dates --- which a panel of 49 large European issuers, all in scope,
does not contain.

The consequence is a restriction on the claim rather than a hedge around it. What the series
identifies is that ESG-lexical content per unit of extracted text rose across 2018--2024 while
tone did not. Whether the added text reflects changed conduct, compliance with a new reporting
regime, or unchanged conduct newly described in a prescribed vocabulary is not identified here,
and the second and third of those are the ones an informed reader should regard as most likely
to be doing work.

There is a sharper consequence, and it belongs to \S\ref{sec:artefact}'s argument rather than
to this section's. Under the mandate reading the index does not merely fail to observe conduct.
$\ESGSI = \Z(\SEN) - \Z(\SUS)$ falls when substance rises, so a firm that adds a mandated
taxonomy table is scored as less likely to be greenwashing --- by construction, for having
complied. An instrument whose washing score is mechanically reduced by regulatory compliance is
not measuring the discrepancy between talk and action; over a window of tightening mandates it
is measuring the volume of prescribed talk. That is a property of the architecture we inherit,
not of our corpus, and it applies to any application of this index family across a period of
regulatory change, the original's included. The series is accordingly better read as a further
demonstration of \S\ref{sec:artefact}'s point than as a finding about European firms.

\subsection{Sensitivity of the series to the substance specification}

The direction and the significance of the trend hold under all three substance specifications
of \S\ref{sec:index}: the yearly means decline and the slope is negative at conventional levels
under $\susdens$, $\sustfidf$ and $\suslag$ alike.

The shape does not hold uniformly, and we state the difference rather than smooth it. The
yearly series is strictly monotone under $\susdens$ and under $\sustfidf$. It is \emph{not}
monotone under $\suslag$: 2024 rebounds to $-0.307$, above the 2023 value of $-0.371$
(Table~\ref{tab:temporal-series}), so that specification yields a series that is monotone over
2018--2023 and reverses in the final year. Monotonicity is therefore a claim about our
specification and about the length-normalised TF-IDF variant, and we make it only for those
two; what generalises across all three is the sign and the significance of the slope.

This precision is required rather than fastidious. \S\ref{sec:robustness} has just established
that the three specifications disagree on 28.0\% of the labels and that the choice between them
is not settled by anything the original discloses. A temporal section that then asserted
uniform behaviour across those same specifications would be claiming a stability that the
preceding section demonstrably refutes. The rebound is also the one place in the window where
the two series separate visibly, and it is consistent with \S\ref{sec:robustness}'s diagnosis:
$\suslag$ correlates $0.9733$ with $\BREADTH$, so a year in which coverage broadens without a
commensurate rise in amount moves it differently from a density measure.

\subsection{Sensitivity of the series to the vocabulary}
\label{sec:temporal-vocabulary}

A dictionary index invites one further objection of the same family as the composition ones:
that the decline is a property of the word list rather than of the reports. The check is to
move the boundary of the list and recompute. The audit of \S\ref{sec:lexicon} supplies the
material --- 147 ESG terms present in the extracted text and absent from our vocabulary ---
and each candidate was measured across all 343 documents for frequency and dispersion before
admission, taking the vocabulary from 154 entries to 289.

Under the 289-entry vocabulary the yearly series is essentially unchanged. The slope moves
from $-0.219$ to $-0.194$ per year, each significant at $p < 10^{-7}$, and restoring the 22
entries the same exercise had set aside as sector-specific leaves it at $-0.193$. Both figures
are computed on the expansion's extraction with the corpus held fixed
(\S\ref{sec:robustness-lexicon}), which is why the base slope here is $-0.219$ rather than the
$-0.205$ of Table~\ref{tab:temporal-series}: the two are not differenced, and only the
proportional movement between them is read. The series does turn on where the boundary is
drawn, and this section contains both directions. Moving it outward, by the 135 audited
entries, changes the slope by about a ninth; moving it inward, by striking the 34-term
reporting-frameworks block (Table~\ref{tab:temporal-frameworks}), changes it by about a fifth,
and the two are of the same order rather than of different ones. What the outward check bounds
is the sensitivity of the series to a defensible extension of this vocabulary, built from these
documents and admitted on dispersion measured over them. It bounds nothing about a vocabulary
built on other principles.

This is not a rebuttal of \S\ref{sec:temporal-mandate}, and the two must not be run together.
Vocabulary insensitivity is a statement about our discretion: had we drawn the boundary of this
word list where the audit supports drawing it instead, the trend would be what it is, so the
result is not an artefact of that choice.
The mandate confound is a statement about chronology: some entries name regimes that did not
exist in 2018, and could not have appeared in an early report whatever the firm was doing.
Enlarging the vocabulary with further contemporary ESG terms does nothing about the second,
and there is reason to think it aggravates it. The candidates were harvested from FY2023
reports (\S\ref{sec:lexicon}), so the added entries are drawn from language at the end of the
window; every framework name in the 154-term list is carried forward untouched; and the
enlarged list is not independent of the corpus it is tested on. The bound on the definitional
component is still the one in Table~\ref{tab:temporal-frameworks} --- $-0.155$ per year with
the reporting-frameworks block struck out --- and the vocabulary check neither tightens nor
loosens it. Robustness to the lexicon and identification of what the lexicon counts are
different questions. The first is bounded here, in one direction, for a class of terms selected
on this corpus for the property that makes them least able to move a density ranking; the
second is not settled anywhere in this paper.

\subsection{What the series licenses}

Both components enter the index after score normalisation across the sample
(\S\ref{sec:artefact}'s object, not \S\ref{sec:measurement}'s), and one consequence must be
explicit before any of these numbers is quoted elsewhere. The index has no cardinal meaning.
A mean of $-0.674$ is not a quantity of ESG-washing but a position within the distribution of
this corpus. Levels are not comparable between corpora: no level reported here is comparable
with any figure in \citet{lagasio2024}, or in any other application of the same architecture,
whatever the score-normalisation rule. This holds even though $\susdens$ is the one substance
specification whose raw value is corpus-independent, because the score normalisation
reintroduces sample dependence downstream, at index assembly. The slope inherits the same
property: $-0.205$ per year is in pooled standard deviations of this sample per year, so it is
no more transferable than the levels are, and a slope estimated on another corpus would not be
commensurable with it. Only the within-sample ordering and the direction of change over time
are interpretable.

Within those bounds, the result is this. On a balanced panel of 49 continuously reporting
European firms over 2018--2024, mean tone-minus-substance declines strictly, and the decline is
carried by the substance component, with no detectable movement in tone. It is not a
composition artefact: the panel is balanced by construction, and the decline survives in a
subsample where the corpus's country and document-type confound does not exist. It is not a
length or an extraction-yield artefact, both of which would bias the measure the other way. Its
sign and significance are invariant to the substance specification, though its monotonicity is
not, and they survive at roughly three quarters of magnitude when every reporting-framework
term is struck from the vocabulary. What it is not shown to be is evidence about conduct, and
the reporting mandates that arrived over the same window supply a specific mechanism by which
the index would move as it does with conduct held fixed (\S\ref{sec:temporal-mandate}) --- a
mechanism whose purely definitional part we can bound but whose remainder no design available
to us can separate from calendar time. That is what the corrected instrument yields
over time: a demonstration of consequences, not the paper's claim.
